\section{Limitations and Conclusion}

This project has several limitations. First, it is simulated tool-call parsing, not real phone automation. The results say whether a model can emit a structured call from a dataset example; they do not measure latency, safety, UI grounding, permission handling, or execution on Android. Second, BabyA and BabyB are compared with each other, but there is no randomly initialized scratch control in the final experiment. Therefore the report can say that BabyB's stronger pretraining appears helpful relative to BabyA, but it cannot prove pretraining is the only cause of the difference. Third, the 128-token context forces compressed prompts and tool signatures, which may make the task harder than it would be for a larger model or a classifier-style architecture.

The final conclusion is deliberately modest. BabyLM-style pretraining helps tiny models move toward structured mobile action parsing: BabyB learns more parseable outputs and more correct function names than BabyA under identical fine-tuning. But BabyA/B are not reliable local phone controllers in this setup. The DSL diagnostic shows that serialization format can create a tiny exact-match signal, while FunctionGemma shows that a small model designed for function calling can solve many more examples on the same evaluation set. BabyActionLM therefore demonstrates both the promise and the current limitation of very small neural language models for local tool-call parsing.
